\section*{Rahmen}
\addcontentsline{toc}{section}{Rahmen}
\markright{Rahmen}
\label{sec:rahmen}

\subsection*{Frage}

Lohnt der Einstieg in einen der sechs Titel zum heutigen Kurs, und bis zu welchem
Kurs? Die Titel stammen aus dem kombinierten Filter vom 19.~September 2026, der
die \FilterTitel{} von Finanzzeitschriften zwischen Juni und September 2026 zum Kauf
empfohlenen Titel nach Wachstums- und Sicherheitskriterien bewertet hat. Sie
erf\"ullen beide Kriterienb\"undel zugleich (\enquote{Dual Champion}). Gew\"ahlt
wurden die beiden punktbesten Titel, dazu je der beste Titel aus einer anderen
Branche und ein Branchenpartner, damit jeder Titel gegen einen vergleichbaren
gelesen werden kann.

\bk{Die Auswahl hat einen eingebauten Fehler, und er ist der Grund f\"ur die
Paarbildung: Das Wachstumskriterium des Filters misst die Umsatzver\"anderung
gegen\"uber dem Vorjahr. In einer zyklischen Branche belohnt es den Gipfel des
Zyklus. Sechs der acht punktbesten Titel waren \"Ol- und Gaswerte. Jede Rechnung
dieses Dokuments stellt deshalb die Frage, an welcher Stelle seiner eigenen
Geschichte ein Titel heute steht.}

\subsection*{Umfang}

Vertiefter Filter, kein Vollbericht. Je Titel: Eigent\"umer und Directors'
Dealings der letzten zw\"olf Monate (Teil~1.2 des Musters), Wettbewerbsposition
(1.6), Mehrjahresreihe des freien Mittelzuflusses (2), Kursverlauf (4) und
Votum mit Schwellenkurs, erforderlichem Wachstum, erforderlichem Endwert und
Multiple im Kontext (7). Nicht bearbeitet: Organe (1.1), Analystenkonsens (1.4),
strategische Ma{\ss}nahmen (3) und Handlungsoptionen (5). Ein Votum aus diesem
Filter ist deshalb ein Votum zur Bewertung, nicht zur Strategie.

\subsection*{Methode in einem Absatz}
\label{sec:methode}

Die Kapitalflussrechnungen wurden aus den gedruckten Seiten der Jahresberichte
gelesen, jede Seite \"uber ihre gedruckte Seitenzahl. Ein lokales Sprachmodell
benannte die Zeilen, Code las die Werte; \GeprueftAnzahl{} Werte wurden gegen den Text
der zitierten Seite gepr\"uft, keiner fehlt. Wo zwei Berichte dasselbe Jahr
drucken, m\"ussen die Lesungen \"ubereinstimmen; Abweichungen sind entweder als
Anpassung des sp\"ateren Berichts erkl\"art oder brechen den Bau ab. F\"ur die
vier SEC-Emittenten wurde jede Reihe zus\"atzlich gegen die XBRL-Daten der SEC
gehalten; Jahre ohne lesbaren Abschluss sind aus XBRL gef\"ullt und in den
Tabellen als Sekund\"arquelle ausgewiesen.

\annahme{Freier Mittelzufluss = operativer Mittelzufluss abz\"uglich
Investitionen in Sachanlagen, Schiffe, \"Ol- und Gasverm\"ogen und immaterielle
Werte, abz\"uglich der Tilgung von Leasingverbindlichkeiten (IFRS-Bilanzierer)
und der Vorzugsdividende (Occidental). Finanzanlagen, Unternehmenserwerbe und
Verkaufserl\"ose bleiben au{\ss}en vor. Bei Reedereien z\"ahlen auch Schiffsk\"aufe
zu den Investitionen, weil die Flotte altert und ersetzt werden muss.}

\annahme{Anlegerrechnung je Titel, unternehmensweit: Einstieg zum heutigen
B\"orsenwert, j\"ahrlicher Zufluss in H\"ohe des freien Mittelzuflusses,
konstant gehalten, Endwert in H\"ohe des Buchwerts des Stammeigenkapitals 2025.
Die Rechnung ist vor Steuern des Anlegers und prognostiziert nichts: Sie h\"alt
das Niveau fest und fragt, welcher Preis dazu passt und welches Wachstum der
heutige Preis bereits enth\"alt.}

\annahme{H\"urde \Pz{\Huerde}: die annualisierte Zehnjahresrendite des MSCI
World einschlie{\ss}lich Dividenden\QW{MSCI World Index (USD), Factsheet}{quelle:msci}{19.09.2026},
also die Rendite der Alternative, die ein Anleger ohne Einzeltitelwahl h\"atte.
Sie ist hoch, weil das Jahrzehnt au{\ss}ergew\"ohnlich war; jede Schwellentabelle
zeigt deshalb eine Spalte mit \Pz{\HuerdeLang}, der Rendite desselben Index seit
Ende 1987 (gleiche Quelle).}

\bk{Drei Szenarien je Titel, alle aus der eigenen Reihe des Unternehmens: der
Median (Basis), das 25.~Perzentil (pessimistisch) und das Gesch\"aftsjahr 2025.
Keines ist eine Prognose. Liegt 2025 hoch in der eigenen Reihe, ist die Basis
das ehrlichere Szenario; liegt es tief, das Gesch\"aftsjahr.}
